% Generated draft: requires mathematical and source review.
\documentclass[11pt]{article}
\usepackage[margin=25mm]{geometry}
\usepackage{fontspec}
\setmainfont{Latin Modern Roman}
\usepackage{amsmath,amssymb,mathtools,unicode-math}
\setmathfont{Latin Modern Math}
\usepackage{enumitem,needspace,xurl,hyperref}
\setlength{\parindent}{0pt}
\setlength{\parskip}{5pt}
\setlength{\emergencystretch}{3em}
\newcommand{\N}{\mathbb{N}}
\newcommand{\Z}{\mathbb{Z}}
\newcommand{\Q}{\mathbb{Q}}
\newcommand{\R}{\mathbb{R}}
\newcommand{\C}{\mathbb{C}}
\newcommand{\F}{\mathbb{F}}
\newcommand{\A}{\mathbb{A}}
\newcommand{\PP}{\mathbb{P}}
\newcommand{\E}{\mathbb{E}}
\newcommand{\eps}{\varepsilon}
\newcommand{\dd}{\,\mathrm{d}}
\begin{document}
\section*{UnsolvedMath ID: 30001249}
\textbf{Source problem code:} catalogue entry
\subsection{Definitions and mathematical statement}
For an integer \(n\geq 1\) and \(p\in[0,1]\), the standard binomial random graph \(G(n,p)\) has vertex set \(\{1,\ldots,n\}\), with each unordered pair of distinct vertices included as an edge independently with probability \(p\). The source does not specify whether \(p\) is fixed or may depend on \(n\), or which range of \(p\) is intended.

Let \(H\) be a fixed graph, independent of \(n\), and let \(X_H(n,p)\) denote the number of copies of \(H\) in \(G(n,p)\). This notation is provisional: the supplied record does not specify the graph class, whether copies are induced, or whether copies are counted as embeddings or as unlabelled subgraphs. These choices must be resolved before \(X_H(n,p)\) is a uniquely defined statistic.

\textbf{Source question.} For every fixed graph \(H\), do analogous modular-equidistribution and local-limit theorems hold for \(X_H(n,p)\)?

\textbf{Unresolved qualifications.} The word ``analogous'' refers to preceding results that are not included in the supplied record. A precise modular-equidistribution assertion requires a modulus or a range of moduli, the residue classes under consideration, a limiting regime, and an error bound or convergence criterion. A precise local-limit assertion requires a normalization, a candidate limiting density, the relevant lattice and its span, and the range over which probability estimates are to hold. None of these is specified here. Degenerate choices of \(H\) or \(p\), and any resulting restrictions on attainable counts or residues, also require explicit treatment. Thus the supplied question is preserved as a research target, but it does not determine a unique fully quantified conjecture.
\subsection{Short English statement}
For each fixed graph \(H\), ask whether its copy count in \(G(n,p)\) satisfies suitable modular-equidistribution and local-limit theorems. The supplied record leaves the comparison theorems, counting convention, parameter ranges, and limiting assertions unspecified, so a fully precise conjecture requires further source context.
\subsection{Sources}
\begin{itemize}
\item Supplied local UnsolvedMath registry record: permanent ID \(30001249\), source problem code \texttt{catalogue entry}, title \emph{Modular and Local Limit Theorems for Subgraph Counts}; Oberwolfach Reports Open Problems collection, proposed year 2009. No individual proposer is identified.
\item Supplied report citation: \emph{Combinatorics and Probability} (2009), Oberwolfach Reports, DOI 10.4171/owr/2009/23, pp.\ 1225--1302. The record describes these reports as documenting mathematical workshops and their open-problem discussions. \url{https://doi.org/10.4171/owr/2009/23}
\item Supplied literature source: ``Sah-Sawhney local limit theorems.'' \url{https://arxiv.org/abs/2006.11369}
\item Supplied literature source: ``joint subgraph-count local limits.'' \url{https://arxiv.org/abs/2412.09535}
\item Literature assessment dated exactly \texttt{2026-08-21}: the supplied curation labels the problem partially solved and reports that later local-limit work proves broad analogues for connected subgraph counts and related joint distributions, while the conjectured full modular statement for every fixed \(H\) is still not a single general theorem. This is supplied dated curation derived from the OWR record, not an independently verified literature assessment; its time-sensitive status and cited claims require independent verification.
\item The registry labels the statement \texttt{corrected\_verified}, explaining that this denotes meaning-preserving semantic, typographic, and LaTeX cleanup with no unsupported source-dependent content added. This supplied label does not resolve the formulation gaps identified above.
\end{itemize}
\end{document}
